\section{Mecanismo interno de las RNN}

\subsection{Intuición de las RNN}

El docente construye la intuición de las RNN mediante un fragmento de pseudocódigo en Python:

\begin{figure}[H]
\centering
\includegraphics[width=0.85\textwidth]{slides-images/slide-20.jpg}
\caption{Intuición del pseudocódigo de una RNN: se recorre cada palabra de la oración, se actualiza el estado oculto y se genera una predicción\protect\footnotemark}
\end{figure}
\footnotetext{Diapositivas, página 20.}

\begin{lstlisting}[caption={Ejemplo de pseudocódigo de una RNN}]
my_rnn = RNN()
hidden_state = [0, 0, 0, 0]

sentence = ["I", "love", "recurrent", "neural"]

for word in sentence:
    prediction, hidden_state = my_rnn(word, hidden_state)

next_word_prediction = prediction
# >>> "networks!"
\end{lstlisting}

Este pseudocódigo muestra con claridad el flujo de trabajo de una RNN:
\begin{enumerate}
    \item Se inicializa el estado oculto como un vector de ceros
    \item Para cada elemento de la secuencia, se invoca la RNN con la entrada actual y el estado oculto del paso anterior
    \item Cada invocación devuelve una predicción y el estado oculto actualizado
    \item Una vez procesada toda la secuencia, la predicción final es la salida
\end{enumerate}

\subsection{Fórmula matemática de la actualización de estado de una RNN}

\begin{figure}[H]
\centering
\includegraphics[width=0.75\textwidth]{slides-images/slide-22.jpg}
\caption{Mecanismo de actualización de estado y salida de una RNN\protect\footnotemark}
\end{figure}
\footnotetext{Diapositivas, página 22.}

La propagación hacia adelante de una RNN comprende dos pasos clave:

\textbf{Paso 1: actualizar el estado oculto}

$$h_t = \tanh(W_{xh} \cdot x_t + W_{hh} \cdot h_{t-1} + b_h)$$

\begin{itemize}
    \item $W_{xh}$: matriz de pesos que transforma la entrada al espacio oculto
    \item $W_{hh}$: matriz de pesos de la actualización del propio estado oculto
    \item $b_h$: término de sesgo
    \item $\tanh$: función de activación no lineal
\end{itemize}

\textbf{Paso 2: generar la salida}

$$\hat{y}_t = W_{hy} \cdot h_t + b_y$$

\begin{itemize}
    \item $W_{hy}$: matriz de pesos que transforma el estado oculto en la salida
    \item $b_y$: sesgo de la salida
\end{itemize}

\begin{importantbox}{Compartición de pesos: la característica clave de las RNN}
En una RNN, las matrices de pesos $W_{xh}$, $W_{hh}$ y $W_{hy}$ se comparten en \textbf{todos los pasos de tiempo}. Esto significa que, sin importar la longitud de la secuencia, el número de parámetros no cambia. Esto no solo reduce en gran medida los parámetros del modelo, sino que además permite que la RNN procese secuencias de cualquier longitud.
\end{importantbox}

\subsection{Despliegue del grafo de cómputo de una RNN}

\begin{figure}[H]
\centering
\includegraphics[width=0.85\textwidth]{slides-images/slide-25.jpg}
\caption{Despliegue del grafo de cómputo de una RNN a lo largo del eje temporal\protect\footnotemark}
\end{figure}
\footnotetext{Diapositivas, página 25.}

Al desplegar la RNN a lo largo del eje temporal, se aprecia con claridad:
\begin{itemize}
    \item Cada paso de tiempo recibe la entrada $x_t$, que participa en el cálculo del estado oculto a través de $W_{xh}$
    \item El estado oculto $h_t$ se transmite al siguiente paso de tiempo a través de $W_{hh}$
    \item Cada paso de tiempo produce la predicción de salida $\hat{y}_t$ a través de $W_{hy}$
\end{itemize}

\subsection{Implementación en código de una RNN}

\begin{lstlisting}[caption={Implementación de una capa RNN con TensorFlow}]
class MyRNNCell(tf.keras.layers.Layer):
    def __init__(self, rnn_units):
        super(MyRNNCell, self).__init__()
        self.W_xh = self.add_weight([rnn_units, rnn_units])
        self.W_hh = self.add_weight([rnn_units, rnn_units])
        self.W_hy = self.add_weight([rnn_units, num_outputs])

    def call(self, x, h):
        # Update hidden state
        h = tf.math.tanh(
            tf.linalg.matmul(x, self.W_xh) +
            tf.linalg.matmul(h, self.W_hh)
        )
        # Compute output
        output = tf.linalg.matmul(h, self.W_hy)
        return output, h
\end{lstlisting}

En marcos de trabajo como TensorFlow/Keras, también se puede usar directamente una capa RNN ya encapsulada:

\begin{lstlisting}[caption={Uso de la capa SimpleRNN encapsulada de Keras}]
tf.keras.layers.SimpleRNN(rnn_units)
\end{lstlisting}

\subsection{Resumen de la sección}

La operación central de una RNN puede resumirse en tres pasos: (1) recibir la entrada actual; (2) combinarla con el estado oculto del paso anterior para actualizar el estado oculto actual; (3) producir una salida a partir del estado oculto. La compartición de pesos en todos los pasos de tiempo es la clave que permite a la RNN procesar secuencias de longitud variable.

\newpage
